% !TeX program = lualatex

\documentclass[14pt]{extarticle}
\usepackage[drawstamps=true]{bstunote}
\usepackage[a4paper, left=23mm, right=9mm, top=9mm, bottom=24mm]{geometry}
\usepackage[english, russian]{babel}

% Включаем гиперссылки в оглавлении
\usepackage[colorlinks=true,linkcolor=blue,urlcolor=blue,citecolor=green!70!black,hypertexnames=false]{hyperref}
\usepackage{lipsum}

\usepackage{fontspec}

\defaultfontfeatures{Mapping=tex-text}

% Шрифт в рамках
\newfontface{\GOSTFontFamily}{GOST-type-A} [ % Название шрифта, который будет в рамках
Path = ./fonts/,
Extension = .ttf,
FakeSlant=0.3 	% Используем искусственный наклон
]

% Шрифт с засечками
\setmainfont{PT-Astra-Serif}[
Path = ./fonts/PT-Astra-Serif/,
Extension = .ttf,
UprightFont = *-Regular,
BoldFont = *-Bold,
ItalicFont = *-Italic,
BoldItalicFont = *-BoldItalic]

% Шрифт без засечек
\setsansfont{PT-Astra-Sans}[
Path = ./fonts/PT-Astra-Sans/,
Extension = .ttf,
UprightFont = *-Regular,
BoldFont = *-Bold,
ItalicFont = *-Italic,
BoldItalicFont = *-BoldItalic]

% Моноширинный шрифт
\setmonofont{PT-Mono}[
Path = ./fonts/PT-Mono/,
Extension = .ttf,
BoldFont = *-Bold,
UprightFont = *-Regular,
Scale=0.9
]

% В LuaLaTeX включаем продвинутую микротипографику 
\usepackage[protrusion=true, expansion=true]{microtype}

\StampAuthor{Иванов И.И.}{3.06.26\,}
\StampCheckBy{Осипов О.В.}{3.06.26\,}
\StampNormControl{Бондаренко Т.В.}{11.06.26\,}
\StampApprovedBy{Поляков В.М.}{25.06.26\,}
\StampWorkTheme{Разработка программного обеспечения для прогнозирования погоды на основе анализа спутниковых снимков с использованием методов машинного обучения}
\StampOrganization{БГТУ~им.~В.Г.~Шухова, \makebox{ПВ-222}}
\StampFont{\GOSTFontFamily}

\begin{document}

\newgeometry{left=3cm, right=1.5cm, top=1.5cm, bottom=1.5cm}
\begin{titlepage}
	\centering % Горизонтальное центрирование
	{\small
	\textbf{МИНОБРНАУКИ РОССИИ} \\[0.5em]
	ФЕДЕРАЛЬНОЕ ГОСУДАРСТВЕННОЕ БЮДЖЕТНОЕ ОБРАЗОВАТЕЛЬНОЕ УЧРЕЖДЕНИЕ 
	ВЫСШЕГО ОБРАЗОВАНИЯ \\[0.3em]
	\textbf{<<БЕЛГОРОДСКИЙ ГОСУДАРСТВЕННЫЙ ТЕХНОЛОГИЧЕСКИЙ\\УНИВЕРСИТЕТ~им.~В.Г.~ШУХОВА>>}\\
	\textbf{(БГТУ~им.~В.Г.~Шухова)}}
	\vspace*{\fill}\par % Вертикальный отступ (заполняет пространство)
	% Здесь размещаем содержимое титульного листа
	{\Large\bfseries Название документа\par}
	\vspace{2cm}
	{\normalsize Автор: Иванов Иван Иванович\par}
	\vspace{\fill} % Ещё один вертикальный отступ
	{\normalsize Белгород \the\year \ г. \par}
\end{titlepage}
\restoregeometry
\initpagenumbering
\setcounter{page}{2}

\section*{OПРЕДЕЛЕНИЯ, СOКРAЩЕНИЯ И OБOЗНAЧЕНИЯ}
\par РФ -- Российская Федерация.
\par СВО -- специальная военная операция.
\par ВУЗ -- высшее учебное заведение.
\par API (Application Programming Interface) -- программный интерфейс, обеспечивающий взаимодействие между компонентами программного обеспечения.
\par RGB (red, green, blue) -- цветовой формат машинной графики, основанный на смешении компонентов красного, зелёного и синего цветов.
\par WWW (World Wide Web) -- всемирная паутина, система взаимосвязанных гипертекстовых документов.

% Содержание
\renewcommand{\contentsname}{\rucontentsname}
\tableofcontents

\section*[ВВЕДЕНИЕ]{ВВЕДЕНИЕ}
\lipsum[1-10]
\section{АНАЛИЗ СОСТОЯНИЯ ИССЛЕДОВАНИЙ И РАЗРАБОТОК В ОБЛАСТИ ПРОГНОЗИРОВАНИЯ ПОГОДЫ}
\lipsum[11]
\subsection{Обзор работ в области прогнозирования погоды}
\subsection{Обзор существующих аналогов, методов и технологий разработки программ для прогнозирования погоды}
\subsection{Объект и предмет исследования}
\subsection{Постановка задачи прогнозирования погоды методом машинного обучения}
\subsection{Практическая и теоретическая значимость работы}
\subsection{Используемые методы и технологии разработки ПО}
\subsection{Описание методов прогнозирования погоды с использованием машинного обучения}
%\subsection{Краткие теоретические сведения}

\lipsum[21-30]
\section{РАЗРАБОТКА АРХИТЕКТУРЫ ПРОГРАММНОГО ОБЕСПЕЧЕНИЯ ДЛЯ ПРОГНОЗИРОВАНИЯ ПОГОДЫ}
\lipsum[31-40]
\section{РАЗРАБОТКА АЛГОРИТМОВ И ПРОГРАММНОГО ОБЕСПЕЧЕНИЯ ДЛЯ ПРОГНОЗИРОВАНИЯ ПОГОДЫ}
\lipsum[41-50]
\section*[ЗАКЛЮЧЕНИЕ]{ЗАКЛЮЧЕНИЕ}
\lipsum[51-60]

\newpage
% Оформление списка литературы
\renewcommand{\refname}{\rubibliographyname} % Заголовок
\phantomsection
\addcontentsline{toc}{section}{\rubibliographyname}% Добавим список литературы в оглавление
\begin{thebibliography}{99}
	\bibitem{gost19201}
	\textbf{ГОСТ 19.201-78.} Техническое задание. Требования к содержанию и оформлению.  -- Введ. 1980-01-01. -- Москва : Стандартинформ, 2010. --~4~с.
\end{thebibliography}

% ПРИЛОЖЕНИЯ
\appendix

\section{Блок-схемы основных алгоритмов прогнозирования погоды}
\lipsum[1-10]
\section{Листинги разработанных подпрограмм для прогнозирования погоды}
\lipsum[11-20]
\end{document}